\documentclass[12pt,a4paper]{article}
\usepackage[utf8]{inputenc}
\usepackage[T1]{fontenc}
\usepackage[spanish]{babel}
\usepackage{geometry}
\usepackage{graphicx}
\usepackage{listings}
\usepackage{xcolor}
\usepackage{hyperref}
\usepackage{fancyhdr}
\usepackage{titlesec}
\usepackage{booktabs}
\usepackage{enumitem}
\usepackage{tcolorbox}
\usepackage{float}

\geometry{margin=2.5cm}

% Colores
\definecolor{codeblue}{RGB}{0,80,180}
\definecolor{codegray}{RGB}{100,100,100}
\definecolor{codegreen}{RGB}{0,120,60}
\definecolor{codebg}{RGB}{245,245,250}
\definecolor{errorred}{RGB}{180,30,30}

% Estilo de código
\lstset{
    backgroundcolor=\color{codebg},
    basicstyle=\ttfamily\small,
    breakatwhitespace=false,
    breaklines=true,
    captionpos=b,
    commentstyle=\color{codegreen},
    keywordstyle=\color{codeblue}\bfseries,
    stringstyle=\color{errorred},
    numbers=left,
    numberstyle=\tiny\color{codegray},
    numbersep=8pt,
    showspaces=false,
    showstringspaces=false,
    showtabs=false,
    frame=single,
    rulecolor=\color{codegray},
    tabsize=2,
    xleftmargin=1.5em,
    framexleftmargin=1em
}

\lstdefinelanguage{powershell}{
    morekeywords={foreach, if, else, param, switch, Write-Host, multipass, sudo, bash, systemctl, curl, Start-Sleep},
    sensitive=false,
    morecomment=[l]{\#},
    morestring=[b]",
}

% Encabezado
\pagestyle{fancy}
\fancyhf{}
\fancyhead[L]{\small Clúster Distribuido con Multipass}
\fancyhead[R]{\small \thepage}
\renewcommand{\headrulewidth}{0.4pt}

\hypersetup{
    colorlinks=true,
    linkcolor=codeblue,
    urlcolor=codeblue,
}

\title{
    \vspace{-1cm}
    {\LARGE \textbf{Documentación del Clúster de Sistemas Distribuidos}} \\[0.5em]
    {\large Implementación con Canonical Multipass en Windows}
}
\author{Proyecto de Sistemas Distribuidos}
\date{\today}

\begin{document}

\maketitle
\tableofcontents
\newpage

% ============================================================
\section{Introducción}

Este documento describe la infraestructura y los comandos utilizados para implementar un clúster de \textbf{3 nodos} de sistemas distribuidos utilizando \textbf{Canonical Multipass} sobre Windows. El clúster implementa:

\begin{itemize}
    \item Máquinas virtuales Ubuntu 24.04 LTS con recursos compartidos
    \item Almacenamiento compartido vía \textbf{NFS} (Network File System)
    \item Comunicación inter-nodos con \textbf{SSH sin contraseña}
    \item Motor de contenedores \textbf{Docker} en cada nodo
    \item Cola de tareas distribuida con \textbf{Redis} y \textbf{Python (Flask)}
\end{itemize}

% ============================================================
\section{Arquitectura del Sistema}

\subsection{Topología de Red}

\begin{table}[H]
\centering
\begin{tabular}{@{}llll@{}}
\toprule
\textbf{Nodo} & \textbf{Rol Principal} & \textbf{Servicios} & \textbf{Imagen} \\
\midrule
node1 & Servidor / Productor & Redis, Flask API, NFS Server & Ubuntu 24.04 \\
node2 & Worker / Cliente & Worker Python, NFS Client & Ubuntu 24.04 \\
node3 & Worker / Cliente & Worker Python, NFS Client & Ubuntu 24.04 \\
\bottomrule
\end{tabular}
\caption{Distribución de roles en el clúster}
\end{table}

\subsection{Recursos por Nodo}

Cada nodo se configura con los siguientes recursos:
\begin{itemize}
    \item \textbf{CPU:} 2 núcleos virtuales
    \item \textbf{RAM:} 2 GB
    \item \textbf{Disco:} 10 GB
\end{itemize}

\subsection{Diagrama de Arquitectura}

\begin{tcolorbox}[colback=codebg, colframe=codegray, title=Flujo de Procesamiento Distribuido]
\begin{verbatim}
  Cliente HTTP
       |
       v
  [node1: Flask :5000]  ---LPUSH--->  [node1: Redis :6379]
                                          |           |
                                       BRPOP       BRPOP
                                          |           |
                                          v           v
                                  [node2: Worker] [node3: Worker]
                                          \           /
                                           \         /
                                        [/shared (NFS)]
\end{verbatim}
\end{tcolorbox}

% ============================================================
\section{Estructura de Archivos del Proyecto}

\begin{lstlisting}[language={}, title=Estructura del directorio]
Multipass_Cluster/
  multipass-init.yaml       # Config cloud-init
  setup_cluster.ps1         # Paso 1: Lanzar nodos
  configure_cluster.ps1     # Paso 2: Configurar NFS
  deploy_app.ps1            # Paso 3: Desplegar app
  manage_cluster.ps1        # Utilidad de gestion
  app/
    producer.py             # API Flask (productor)
    worker.py               # Consumidor de tareas
    requirements.txt        # Dependencias Python
\end{lstlisting}

% ============================================================
\section{Paso 1: Inicialización del Clúster (setup\_cluster.ps1)}

\subsection{Cloud-Init (multipass-init.yaml)}

El archivo de cloud-init configura cada nodo automáticamente al momento de su creación:

\begin{lstlisting}[language=bash, title=multipass-init.yaml]
#cloud-config
users:
  - name: distadmin
    sudo: ALL=(ALL) NOPASSWD:ALL
    groups: [sudo, docker]
    shell: /bin/bash
    lock_passwd: false

chpasswd:
  list: |
     distadmin:tera
  expire: False

packages:
  - openssh-server
  - curl, git, build-essential
  - docker.io
  - nfs-kernel-server, nfs-common
  - net-tools, htop
  - python3, python3-pip, python3-venv

runcmd:
  - [ systemctl, enable, ssh ]
  - [ systemctl, enable, docker ]
  - [ mkdir, -p, /shared ]
  - [ chown, -R, "distadmin:distadmin", /shared ]
  - [ chmod, 777, /shared ]
  - [ mkdir, -p, /home/distadmin/.ssh ]
  - [ chmod, 700, /home/distadmin/.ssh ]
\end{lstlisting}

\subsection{Lanzamiento de Nodos}

\begin{lstlisting}[language=powershell, title=Comando de lanzamiento]
# Lanzar cada nodo con los recursos definidos
multipass launch 24.04 --name node1 --cpus 2 `
    --memory 2G --disk 10G --cloud-init multipass-init.yaml
multipass launch 24.04 --name node2 --cpus 2 `
    --memory 2G --disk 10G --cloud-init multipass-init.yaml
multipass launch 24.04 --name node3 --cpus 2 `
    --memory 2G --disk 10G --cloud-init multipass-init.yaml

# Esperar a que cloud-init termine
multipass exec node1 -- cloud-init status --wait
\end{lstlisting}

\subsection{Configuración de /etc/hosts}

Se obtienen las direcciones IP de cada nodo y se registran en \texttt{/etc/hosts} de todos los nodos para permitir la resolución de nombres:

\begin{lstlisting}[language=powershell, title=Registro de hostnames]
# Obtener IP de un nodo
$ipInfo = multipass info node1 --format csv | ConvertFrom-Csv

# Inyectar entradas en /etc/hosts de cada nodo
multipass exec node1 -- sudo bash -c `
    "echo -e '172.x.x.x node1\n172.x.x.x node2\n172.x.x.x node3' >> /etc/hosts"
\end{lstlisting}

\subsection{SSH Sin Contraseña}

Se genera un par de claves RSA compartido y se distribuye a todos los nodos:

\begin{lstlisting}[language=powershell, title=Distribución de claves SSH]
# Generar clave compartida
ssh-keygen -t rsa -b 4096 -f $keyFile -N '""' -q

# Instalar en cada nodo
multipass exec node1 -- bash -c "echo '$privKey' > /home/distadmin/.ssh/id_rsa"
multipass exec node1 -- chmod 600 /home/distadmin/.ssh/id_rsa
multipass exec node1 -- bash -c "echo '$pubKey' >> /home/distadmin/.ssh/authorized_keys"

# Pre-agregar a known_hosts para evitar el prompt inicial
multipass exec node1 -- sudo -u distadmin bash -c `
    "ssh-keyscan -H node2 >> /home/distadmin/.ssh/known_hosts"
\end{lstlisting}

% ============================================================
\section{Paso 2: Configuración de Servicios (configure\_cluster.ps1)}

\subsection{Servidor NFS (node1)}

\begin{lstlisting}[language=powershell, title=Configuración del servidor NFS]
# Exportar el directorio /shared
multipass exec node1 -- sudo bash -c `
    "echo '/shared *(rw,sync,no_subtree_check,no_root_squash)' > /etc/exports"
multipass exec node1 -- sudo exportfs -arv
multipass exec node1 -- sudo systemctl restart nfs-kernel-server
\end{lstlisting}

\subsection{Clientes NFS (node2, node3)}

\begin{lstlisting}[language=powershell, title=Montaje de NFS en clientes]
# Agregar entrada a fstab para persistencia
multipass exec node2 -- sudo bash -c `
    "echo '${nfsServer}:/shared /shared nfs auto,nofail,noatime,nolock,intr,tcp,actimeo=1800 0 0' >> /etc/fstab"
multipass exec node2 -- sudo mount -a
\end{lstlisting}

\subsection{Verificación de Docker}

\begin{lstlisting}[language=powershell, title=Verificación de Docker]
multipass exec node1 -- docker --version
# Salida esperada: Docker version 27.x.x, build ...
\end{lstlisting}

% ============================================================
\section{Paso 3: Despliegue de la Aplicación (deploy\_app.ps1)}

\subsection{Transferencia de Archivos}

\begin{lstlisting}[language=powershell, title=Transferencia vía multipass]
# Crear directorio destino
multipass exec node1 -- sudo mkdir -p /shared/app
multipass exec node1 -- sudo chmod 777 /shared/app

# Transferir archivos al nodo
multipass transfer producer.py "node1:producer.py"
multipass exec node1 -- sudo cp "producer.py" "/shared/app/producer.py"
\end{lstlisting}

\subsection{Instalación de Redis}

\begin{lstlisting}[language=powershell, title=Configuración de Redis]
# Instalar Redis
multipass exec node1 -- sudo apt-get install -y redis-server

# Configurar para escuchar en todas las interfaces
multipass exec node1 -- sudo bash -c "echo 'bind 0.0.0.0' >> /etc/redis/redis.conf"
multipass exec node1 -- sudo bash -c "echo 'protected-mode no' >> /etc/redis/redis.conf"
multipass exec node1 -- sudo systemctl restart redis-server

# Verificar que Redis escucha en el puerto correcto
multipass exec node1 -- sudo ss -tlnp | findstr 6379
# Esperado: LISTEN 0 511 0.0.0.0:6379 0.0.0.0:*
\end{lstlisting}

\subsection{Entorno Virtual de Python}

\begin{lstlisting}[language=powershell, title=Configuración del entorno Python]
multipass exec node1 -- bash -c `
    "python3 -m venv /shared/app/venv && `
     /shared/app/venv/bin/pip install -r /shared/app/requirements.txt"
\end{lstlisting}

\subsection{Servicios systemd}

En lugar de usar \texttt{nohup} (que no persiste en sesiones de multipass exec), se crean servicios \texttt{systemd}:

\begin{lstlisting}[language=bash, title=producer.service (node1)]
[Unit]
Description=Distributed Task Queue Producer
After=network.target

[Service]
Type=simple
User=distadmin
WorkingDirectory=/shared/app
ExecStart=/shared/app/venv/bin/python /shared/app/producer.py
Restart=always
RestartSec=3
StandardOutput=append:/shared/app/producer.log
StandardError=append:/shared/app/producer.log

[Install]
WantedBy=multi-user.target
\end{lstlisting}

\begin{lstlisting}[language=bash, title=worker.service (node2 y node3)]
[Unit]
Description=Distributed Task Queue Worker
After=network.target

[Service]
Type=simple
User=distadmin
WorkingDirectory=/shared/app
ExecStart=/shared/app/venv/bin/python /shared/app/worker.py
Restart=always
RestartSec=3
StandardOutput=append:/shared/app/worker.log
StandardError=append:/shared/app/worker.log

[Install]
WantedBy=multi-user.target
\end{lstlisting}

\begin{lstlisting}[language=powershell, title=Activación de servicios]
# Activar el productor en node1
multipass exec node1 -- sudo systemctl daemon-reload
multipass exec node1 -- sudo systemctl restart producer.service

# Activar workers en node2 y node3
multipass exec node2 -- sudo systemctl daemon-reload
multipass exec node2 -- sudo systemctl restart worker.service
multipass exec node3 -- sudo systemctl daemon-reload
multipass exec node3 -- sudo systemctl restart worker.service
\end{lstlisting}

% ============================================================
\section{Código de la Aplicación}

\subsection{Productor (producer.py)}

\begin{lstlisting}[language=Python, title=producer.py --- API Flask]
from flask import Flask, request, jsonify
import redis, time, socket

app = Flask(__name__)
r = redis.Redis(host='node1', port=6379, db=0)

@app.route('/')
def index():
    return jsonify({
        "status": "online",
        "service": "producer",
        "hostname": socket.gethostname()
    })

@app.route('/enqueue', methods=['GET', 'POST'])
def enqueue():
    task_id = request.args.get('task', f"task_{int(time.time())}")
    r.lpush('task_queue', task_id)
    return jsonify({
        "message": f"Task '{task_id}' enqueued successfully",
        "queue_length": r.llen('task_queue')
    })

if __name__ == '__main__':
    app.run(host='0.0.0.0', port=5000)
\end{lstlisting}

\subsection{Worker (worker.py)}

\begin{lstlisting}[language=Python, title=worker.py --- Consumidor de tareas]
import redis, time, socket, logging

logging.basicConfig(level=logging.INFO,
    format='%(asctime)s [%(levelname)s] (%(hostname)s) %(message)s')
hostname = socket.gethostname()
logger = logging.LoggerAdapter(
    logging.getLogger(__name__), {'hostname': hostname})

r = redis.Redis(host='node1', port=6379, db=0)

def worker():
    logger.info("Worker started, waiting for tasks...")
    while True:
        try:
            task = r.brpop('task_queue', timeout=5)
            if task:
                queue_name, task_id = task
                task_id = task_id.decode('utf-8')
                logger.info(f"Picked up task: {task_id}")
                time.sleep(2)  # Simular procesamiento
                logger.info(f"Finished processing task: {task_id}")
        except redis.ConnectionError:
            logger.error("Could not connect to Redis. Retrying in 5s...")
            time.sleep(5)

if __name__ == '__main__':
    worker()
\end{lstlisting}

% ============================================================
\section{Verificación y Pruebas}

\subsection{Comandos de Verificación}

\begin{lstlisting}[language=powershell, title=Lista completa de verificaciones]
# 1. Estado del cluster
multipass list

# 2. Conectividad de red entre nodos
multipass exec node2 -- ping -c 2 node1
multipass exec node3 -- ping -c 2 node1

# 3. Almacenamiento compartido NFS
multipass exec node1 -- touch /shared/test.txt
multipass exec node3 -- ls /shared/test.txt

# 4. Redis escuchando correctamente
multipass exec node1 -- sudo ss -tlnp | findstr 6379

# 5. API del productor respondiendo
multipass exec node1 -- curl -s http://127.0.0.1:5000/
\end{lstlisting}

\subsection{Prueba de Procesamiento Distribuido}

\begin{lstlisting}[language=powershell, title=Prueba end-to-end]
# Enviar 4 tareas al productor
multipass exec node1 -- curl -s "http://127.0.0.1:5000/enqueue?task=AnalyzeDatasetA"
multipass exec node1 -- curl -s "http://127.0.0.1:5000/enqueue?task=AnalyzeDatasetB"
multipass exec node1 -- curl -s "http://127.0.0.1:5000/enqueue?task=ProcessImageBatch"
multipass exec node1 -- curl -s "http://127.0.0.1:5000/enqueue?task=TrainModelAlpha"

# Esperar y revisar logs de los workers
Start-Sleep -Seconds 10
multipass exec node2 -- cat /shared/app/worker.log
multipass exec node3 -- cat /shared/app/worker.log
\end{lstlisting}

\subsection{Resultados Esperados}

La salida de los logs de los workers deberá mostrar:

\begin{lstlisting}[language={}, title=Ejemplo de salida exitosa]
# worker.log en node2:
2026-05-17 16:04:35 [INFO] (node2) Worker started, waiting for tasks...
2026-05-17 16:05:02 [INFO] (node2) Picked up task: AnalyzeDatasetA
2026-05-17 16:05:04 [INFO] (node2) Finished processing task: AnalyzeDatasetA

# worker.log en node3:
2026-05-17 16:04:40 [INFO] (node3) Worker started, waiting for tasks...
2026-05-17 16:05:02 [INFO] (node3) Picked up task: ProcessImageBatch
2026-05-17 16:05:04 [INFO] (node3) Finished processing task: ProcessImageBatch
\end{lstlisting}

% ============================================================
\section{Gestión del Clúster (manage\_cluster.ps1)}

\begin{lstlisting}[language=powershell, title=Comandos de gestión]
# Ver estado de todos los nodos
.\manage_cluster.ps1 -Action status

# Iniciar todos los nodos
.\manage_cluster.ps1 -Action start

# Detener todos los nodos
.\manage_cluster.ps1 -Action stop

# Reiniciar todos los nodos
.\manage_cluster.ps1 -Action restart

# Abrir shell en un nodo especifico
.\manage_cluster.ps1 -Action shell -Node node2

# Destruir el cluster completo
.\manage_cluster.ps1 -Action destroy
\end{lstlisting}

% ============================================================
\section{Errores Conocidos y Soluciones}

\begin{table}[H]
\centering
\small
\begin{tabular}{@{}p{3cm}p{4.5cm}p{5.5cm}@{}}
\toprule
\textbf{Error} & \textbf{Causa} & \textbf{Solución} \\
\midrule
Variable \texttt{\$var:} no válida en PowerShell &
PowerShell interpreta \texttt{\$var:} como cualificador de ámbito &
Usar sintaxis \texttt{\$\{var\}} cuando la variable precede a dos puntos \\
\midrule
Redis bind \texttt{0.0.0.1} &
Error tipográfico en el script original &
Corregir a \texttt{0.0.0.0} \\
\midrule
\texttt{sed} no modifica redis.conf &
El formato del archivo en Ubuntu 24.04 difiere del patrón buscado &
Agregar \texttt{bind 0.0.0.0} y \texttt{protected-mode no} al final del archivo \\
\midrule
SFTP Permission Denied &
Ruta de destino incorrecta en \texttt{multipass transfer} &
Usar ruta relativa: \texttt{node1:archivo} en lugar de ruta absoluta \\
\midrule
Procesos \texttt{nohup} mueren &
\texttt{multipass exec} termina la sesión bash y sus hijos &
Usar servicios \texttt{systemd} para persistencia de procesos \\
\bottomrule
\end{tabular}
\caption{Resumen de errores encontrados y sus correcciones}
\end{table}

% ============================================================
\section{Conclusiones}

Se implementó exitosamente un clúster de 3 nodos de sistemas distribuidos con las siguientes capacidades:

\begin{enumerate}
    \item \textbf{Virtualización:} Tres máquinas virtuales Ubuntu 24.04 administradas por Multipass.
    \item \textbf{Red:} Resolución de nombres entre nodos vía \texttt{/etc/hosts} y comunicación SSH sin contraseña.
    \item \textbf{Almacenamiento compartido:} Directorio \texttt{/shared} sincronizado vía NFS.
    \item \textbf{Procesamiento distribuido:} Cola de tareas Redis consumida por workers en nodos independientes.
    \item \textbf{Gestión de procesos:} Servicios \texttt{systemd} para la ejecución persistente de la aplicación.
\end{enumerate}

El sistema demostró la distribución efectiva de tareas entre \texttt{node2} y \texttt{node3}, verificando que ambos workers consumieron tareas de la cola Redis en \texttt{node1} de forma independiente y concurrente.

\end{document}
